% Manuscript: Regime-Dependent Performance of Learned Kalman Observers
%             for Early-Warning Detection
%
% After selecting a target journal, replace \documentclass{article} with
% the journal template and remove placeholder settings that the template
% replaces.

\documentclass[11pt]{article}

\usepackage[margin=1in]{geometry}
\usepackage[T1]{fontenc}
\usepackage[utf8]{inputenc}
\usepackage{lmodern}
\usepackage{amsmath,amssymb,amsthm}
\usepackage{mathtools}
\usepackage{graphicx}
\usepackage{booktabs}
\usepackage{enumitem}
\usepackage{flafter}
\usepackage[numbers,sort&compress]{natbib}
\usepackage[colorlinks=true,linkcolor=blue,citecolor=blue,urlcolor=blue]{hyperref}
\usepackage[nameinlink,noabbrev]{cleveref}

\graphicspath{{data/figures/}}

\crefname{section}{Section}{Sections}
\Crefname{section}{Section}{Sections}
\crefname{subsection}{Section}{Sections}
\Crefname{subsection}{Section}{Sections}
\crefname{table}{Table}{Tables}
\Crefname{table}{Table}{Tables}
\crefname{figure}{Figure}{Figures}
\Crefname{figure}{Figure}{Figures}

\newtheorem{theorem}{Theorem}[section]
\newtheorem{proposition}[theorem]{Proposition}
\newtheorem{lemma}[theorem]{Lemma}
\newtheorem{corollary}[theorem]{Corollary}
\theoremstyle{remark}
\newtheorem{remark}[theorem]{Remark}

\title{Regime-Dependent Performance of Learned Kalman Observers \\ for Early-Warning Detection}

\author{}

\date{}

\begin{document}

\maketitle

\begin{abstract}
Early-warning methods for critical transitions are often evaluated as if a
single detector should transfer across dynamical regimes. This assumption is
too strong when the warning signatures differ across bifurcation mechanisms.
We study this issue with learned Kalman observer variants on synthetic fold,
Hopf, and logistic systems, using patient-count sweeps from 100 to 500
trajectories and ten random seeds per setting. The benchmark compares a
binary-cross-entropy Kalman observer with a spectral recurrent variant using
detection time, early-warning AUC, and false-positive rate. At 500 patients,
the spectral recurrent observer increases early-warning AUC on fold and Hopf
systems and matches the BCE observer on logistic dynamics, but the gains do
not extend uniformly to false-positive control or detection-time summaries.
Across patient counts, the relative behavior of the two observers remains
system- and metric-dependent. These results argue against treating
early-warning detection as a single transferable classification problem.
Instead, observer design and model selection should be tied to the dynamical
regime, the available data depth, and the operating cost of false alarms
versus delayed warnings.
\end{abstract}

\noindent\textbf{Keywords:} critical transitions; early-warning signals; Kalman observer; bifurcation detection; spectral regularization; dynamical systems

\section{Introduction}
\label{sec:introduction}

Critical transitions in dynamical systems are often preceded by measurable
changes in the observable time series. The practical question is whether
these changes can be detected early enough to support intervention. This
makes early-warning detection a problem of timing, calibration, and
false-alarm control---not only binary classification.

Classical early-warning indicators based on variance, autocorrelation, and
related statistics have been applied across ecology, climate science, and
engineering \citep{scheffer2009early,dakos2008slowing}. These indicators are
transparent and require no model fitting, but each encodes a fixed functional
form for the warning signature. When the transition mechanism changes, a
statistic that was informative for one bifurcation type may become unreliable
for another.

Learned observers offer an alternative. A Kalman-style observer maintains a
latent state updated by a prediction--correction rule and can learn the
mapping from observations to warning scores from data. Recent work has shown
that such models can extract early-warning information from simulated
systems and, in some cases, transfer across dynamics with matching
normal-form structure \citep{bury2021deep,deb2022machine}. However, a
fold bifurcation, a Hopf bifurcation, and a period-doubling transition
arise from different local mechanisms
\citep{kuznetsov2004elements,strogatz2015nonlinear} and produce different
pre-transition signatures. There is no reason to expect uniform transfer
across all of them.

This paper studies the transfer question directly. We compare two learned
Kalman observer variants---a binary-cross-entropy observer (Kalman-BCE) and
a spectral recurrent variant (Kalman-LSTM-Spec)---across synthetic fold,
Hopf, and logistic systems. The benchmark varies the training-set size from
100 to 500 trajectories and evaluates each setting across ten random seeds.
Performance is measured by three metrics: detection time, early-warning AUC,
and false-positive rate. This design separates ranking quality from
operational calibration.

The main finding is regime dependence. At 500 trajectories, the spectral
recurrent observer achieves higher early-warning AUC on fold
($0.903 \pm 0.059$ vs.\ $0.827 \pm 0.235$) and Hopf
($0.823 \pm 0.243$ vs.\ $0.617 \pm 0.177$), while both methods reach
perfect AUC on the logistic system. However, the BCE observer produces
earlier detections on fold (detection time $109.5 \pm 28.6$ vs.\
$72.7 \pm 11.0$ steps) and lower false-positive rates on fold and logistic.
The relative ranking also shifts with training-set size: LSTM-Spec benefits
more from additional data on fold and Hopf, while logistic performance
saturates by 200 trajectories for both methods. No single observer dominates
across all three systems, metrics, and data sizes.

\subsection{Contributions}
\label{sec:contributions}

We make four contributions.

\begin{enumerate}[leftmargin=*]
  \item A controlled benchmark for learned Kalman observers across fold,
        Hopf, and logistic bifurcation systems, with deterministic seed
        generation and batch selection for reproducibility.
  \item Evaluation across five training-set sizes (100--500 trajectories)
        with ten seeds per setting, providing cross-seed variance estimates
        for every reported metric.
  \item A multi-metric comparison using detection time, early-warning AUC,
        and false-positive rate, exposing trade-offs that a single aggregate
        score would hide.
  \item Evidence that observer performance is system- and metric-dependent:
        the spectral recurrent observer improves discrimination on fold and
        Hopf but does not dominate the BCE observer across all systems,
        data sizes, and metrics.
\end{enumerate}

These contributions support a bounded claim: learned observer design for
early-warning detection should be evaluated as a regime-dependent modeling
problem, not as a search for a detector that transfers uniformly across
bifurcation mechanisms.

\section{Related Work}
\label{sec:related}

% If the target venue uses strict IMRAD formatting, this section can be
% merged into Section~\ref{sec:introduction}.

\subsection{Classical Early-Warning Signals}
\label{sec:related:classical}

% Variance, autocorrelation, lag-based indicators, critical slowing down.
% Role: establishes the non-learned baseline context; motivates why simple
% indicators can remain competitive; explains why learned methods should be
% compared against strong simple baselines.

\subsection{Observer and Kalman-Based Approaches}
\label{sec:related:observer}

% Kalman filtering, observer dynamics, and learned observer variants.
% Role: motivates the use of a Kalman observer; explains why stability and
% filtering behavior matter; prepares the reader for the learned observer design.

\subsection{Learning Across Dynamical Regimes}
\label{sec:related:regimes}

% Why methods may not transfer uniformly across regimes. Role: supports the
% honest framing that method performance can be conditional; prepares reviewers
% for a trade-off result rather than a universal-best claim.

\section{Problem Setting and Benchmark Definition}
\label{sec:problem}

We evaluate early-warning detectors on simulated trajectory data from three
dynamical systems. Each detector receives a trajectory and produces a score
sequence. The score at time \(t\), denoted \(p_{i,t}\), is interpreted as the
detector's warning score for trajectory \(i\). Higher values indicate stronger
evidence for an approaching transition.

The benchmark is defined to compare observer behavior across regimes rather
than to select a single global winner. All reported comparisons use the same
systems, patient-count settings, seed order, thresholding rule, and evaluation
metrics.

\subsection{Dynamical Systems}
\label{sec:systems}

The benchmark contains three synthetic systems.

\begin{description}[leftmargin=*]
  \item[Fold.]
  The fold system represents a collapse-like transition in which a stable state
  is lost as the control condition changes. It tests whether an observer can
  detect gradual loss of stability without relying on oscillatory structure.

  \item[Hopf.]
  The Hopf system represents an oscillatory instability. It tests whether an
  observer can use temporal structure associated with changing stability and
  emerging oscillatory behavior.

  \item[Logistic.]
  The logistic system represents a discrete-time transition with
  period-doubling behavior. It tests whether an observer calibrated on
  continuous warning scores also behaves well for map-like dynamics.
\end{description}

These systems are not treated as interchangeable datasets. They are distinct
test cases for the same observer family.

\subsection{Trajectory Data and Patient Counts}
\label{sec:data}

For each system, the benchmark contains signal trajectories and null
trajectories. Signal trajectories contain a transition with a bifurcation time
\(\tau_i\). Null trajectories provide non-transition examples for
false-positive evaluation.

The patient count denotes the number of simulated trajectories used in a
benchmark condition. The finalized patient-count settings are
\[
  n \in \{100, 200, 300, 400, 500\}.
\]
For each patient count, system, and observer variant, results are aggregated
over the following ten seeds:
\[
  101, 202, 303, 404, 505, 606, 707, 808, 909, 1010.
\]

The finalized reporting pipeline compares two learned observer variants:
\(\text{Kalman-BCE}\) and \(\text{Kalman-LSTM-Spec}\). Each method produces
test scores \(p_{i,t} \in [0,1]\), which are used for threshold selection,
detection-time computation, early-warning AUC, and false-positive evaluation.

\subsection{Batch Selection}
\label{sec:batch}

The experiment directory can contain multiple batches for the same patient
count, including partial reruns. To avoid manual selection, the analysis uses a
deterministic batch-selection rule. For each patient count \(n\), let
\(\mathcal{B}_n\) be the set of available batches. The selected batch is the
batch that maximizes, in order:

\begin{enumerate}[leftmargin=*]
  \item the number of unique \((\text{system}, \text{method}, \text{seed})\)
        triples,
  \item the number of systems represented,
  \item the total number of result records,
  \item and the timestamp, if the preceding quantities are tied.
\end{enumerate}

This rule keeps the most complete batch for each patient count and excludes
partial reruns from the main figures and tables.

\subsection{Evaluation Metrics}
\label{sec:metrics}

The benchmark reports three metrics: detection time, early-warning AUC, and
false-positive rate. The metrics measure different properties of the same score
sequence and should not be collapsed into a single ordering.

\subsubsection{Detection Time}
\label{sec:metrics:dt}

For a positive trajectory \(i\) with bifurcation time \(\tau_i\), define the
first pre-bifurcation alert time as
\[
  a_i = \min \{t : 0 \leq t < \tau_i,\; p_{i,t} \geq \theta \},
\]
where \(\theta\) is the decision threshold. If such an alert exists, the
trajectory-level detection time is
\[
  DT_i = \tau_i - a_i.
\]
The reported detection time is the mean of \(DT_i\) over positive trajectories
with at least one pre-bifurcation alert. Trajectories without a pre-bifurcation
alert are omitted from this mean rather than counted as zero-lead detections.

\subsubsection{Early-Warning AUC}
\label{sec:metrics:auc}

For each positive trajectory, the early-warning score is the maximum model
score in a window before the bifurcation:
\[
  e_i = \max \{p_{i,t} : \max(0,\tau_i - 50) \leq t < \tau_i - 5\}.
\]
For each null trajectory \(j\) with sequence length \(T_j\), the corresponding
score is computed over the terminal part of the sequence:
\[
  n_j = \max \{p_{j,t} : \max(0,T_j - 50) \leq t < T_j - 5\}.
\]
The early-warning AUC is the ROC AUC computed from the pooled signal scores
\(\{e_i\}\) and null scores \(\{n_j\}\), with labels one for signal and zero
for null.

\subsubsection{False Positive Rate}
\label{sec:metrics:fpr}

False-positive rate is computed on null trajectories as a per-time-step alert
rate:
\[
  FPR =
  \frac{\sum_j \sum_{t=0}^{T_j-1} \mathbf{1}[p_{j,t} \geq \theta]}
       {\sum_j T_j}.
\]
This quantity measures the fraction of null time steps that would trigger an
alert. It is not a per-trajectory false-positive probability.

\subsubsection{Threshold Selection}
\label{sec:metrics:threshold}

Thresholds are selected on the validation split before test evaluation. For a
positive validation trajectory, scores in the window
\[
  \max(0,\tau_i - 60) \leq t < \tau_i
\]
are labeled positive. Scores from the early pre-bifurcation region
\[
  0 \leq t < \max(\max(0,\tau_i - 120) - 1, 0)
\]
are labeled negative. If null validation scores are supplied, their first 120
time steps are also labeled negative. In the finalized benchmark runs reported
here, threshold selection for the two learned observer variants uses the signal
validation trajectories.

The threshold is selected from the validation ROC curve by maximizing Youden's
statistic \citep{youden1950index}, \(J = TPR - FPR\). If the sensitivity at that threshold is below
0.80, the selected threshold is the one whose sensitivity is closest to 0.80.
If the validation labels contain only one class, or if the validation scores are
effectively constant, the fallback threshold is 0.5.

\section{Methods}
\label{sec:methods}

\subsection{Learned Kalman Observer}
\label{sec:methods:observer}

The observer maintains a latent state \(z_t \in \mathbb{R}^d\) and updates it
at each time step using a prediction--correction rule. Given observation
\(y_t \in \mathbb{R}^D\) at time~\(t\):
\begin{align}
  \hat{z}_t &= A\, z_{t-1}, \label{eq:predict} \\
  z_t &= \hat{z}_t + K\bigl(y_t - C\,\hat{z}_t\bigr), \label{eq:correct}
\end{align}
where \(A \in \mathbb{R}^{d \times d}\) is the state transition matrix,
\(C \in \mathbb{R}^{D \times d}\) maps latent states to predicted observations,
and \(K \in \mathbb{R}^{d \times D}\) is the observer gain. All three matrices
are learned. \(A\) is initialized as \(0.95\,I_d\); entries of \(C\) and \(K\)
are drawn i.i.d.\ from \(\mathcal{N}(0, 0.1^2)\). The latent state starts at
zero.

At each step, a scoring head maps the latent state to a scalar logit
\(\ell_t\). The warning score is \(p_t = \sigma(\ell_t)\), where \(\sigma\)
denotes the sigmoid function. This score is used for threshold selection and
evaluation as defined in Section~\ref{sec:metrics}.

\subsection{Observer Variants}
\label{sec:methods:variants}

Both reported variants share the Kalman core above with latent dimension
\(d = 4\). They differ in the scoring head and, for one variant, in the
transition dynamics.

\begin{description}[leftmargin=*]
  \item[Kalman-BCE.]
  A feedforward head maps the latent state to the logit through layer
  normalization, a linear projection from \(d\) to \(\lfloor d/2 \rfloor\),
  GELU activation, and a linear projection to a scalar. The transition matrix
  \(A\) is fixed across time steps.

  \item[Kalman-LSTM-Spec.]
  An LSTM cell \citep{hochreiter1997long} with hidden dimension~8 receives the latent state at each step.
  Its hidden state \(h_t\) is passed through the same two-layer head structure
  to produce the logit \(\ell_t\). The sigmoid output
  \(\alpha_t = \sigma(\ell_t)\) gates the transition matrix:
  \[
    A_t = (1 - \alpha_t)\,A + \alpha_t\,I_d,
  \]
  replacing \(A\) in~\eqref{eq:predict}. When \(\alpha_t\) is large, the
  transition is pulled toward the identity, attenuating the learned dynamics.
  Training adds a spectral-radius penalty described below.
\end{description}

These two variants are the main comparison in all benchmark tables and figures.

\subsection{Training Objective and Calibration}
\label{sec:methods:training}

The training target for a positive trajectory with bifurcation time \(\tau\) is
\[
  y_t^{*} =
  \begin{cases}
    \exp\!\bigl(-\bigl((\tau - t)/\sigma\bigr)^2\bigr), & t \leq \tau, \\[4pt]
    0, & t > \tau,
  \end{cases}
\]
with \(\sigma = 60\). Null trajectories receive zero targets at all time steps.
The loss is the mean binary cross-entropy between the model logits and
\(y_t^{*}\), averaged over valid (non-padded) time steps.

For Kalman-LSTM-Spec, a spectral-radius penalty is added. The closed-loop
matrix \(M = A - KCA\) is formed from the current learned parameters, and its
spectral radius \(\rho(M)\) is estimated via ten power iterations. The penalty
is \(\lambda\,\max\bigl(\rho(M) - 0.95,\; 0\bigr)\) with \(\lambda = 0.01\).

Models are optimized with AdamW \citep{loshchilov2019decoupled} (learning rate \(10^{-3}\), weight decay
\(10^{-5}\)) and a cosine-annealing schedule \citep{loshchilov2017sgdr} over 50~epochs. Gradients are
clipped at unit norm. Early stopping with patience~5 restores the checkpoint
with the lowest validation loss. Each run seeds NumPy, PyTorch, and CUDA to the
run-level seed before initialization.

The data are split 60/20/20 into training, validation, and test sets. Threshold
selection on the validation split follows Section~\ref{sec:metrics:threshold}.
No test data enter threshold or model selection.

\subsection{Implementation and Reproducibility}
\label{sec:methods:implementation}

The observer and training loop are implemented in PyTorch \citep{paszke2019pytorch}. Experiment
configurations use composable YAML files with separate templates for data
generation, model architecture, and training hyperparameters; patient-count
overrides are the only difference between run configurations.

The benchmark runner executes all \((\text{system},\, \text{method},\,
\text{seed})\) combinations for a given patient count in a single invocation.
Raw outputs are written to \texttt{analysis/experiment\_data/}, organized by
patient count and batch timestamp. The visualization pipeline reads the selected
batches (Section~\ref{sec:batch}) and produces the tables and figures reported
here.

\section{Experimental Design}
\label{sec:design}

This section describes the comparisons reported in the main results. Evaluation
metrics, threshold selection, and batch selection are defined in
Section~\ref{sec:problem}.

\subsection{Main Benchmark}
\label{sec:design:main}

The primary comparison evaluates Kalman-BCE and Kalman-LSTM-Spec across all
three systems at five patient counts (100, 200, 300, 400, 500), each with ten
seeds. This produces 60~runs per patient count:
\(3\;\text{systems} \times 2\;\text{methods} \times 10\;\text{seeds}\).

Table~\ref{tab:benchmark_500} reports the 500-patient results: mean and
standard deviation of detection time, early-warning AUC, and false-positive rate
across seeds, listed per system and method. Table~\ref{tab:benchmark_sweep}
extends this to all five patient counts.

Figure~\ref{fig:patient_sweep} plots detection time and early-warning AUC as a
function of patient count, with separate panels per system. This tests whether
performance changes with training set size and whether the two variants respond
differently to that change.

\subsection{Trajectory-Level Behavior}
\label{sec:design:trajectory}

Figure~\ref{fig:trajectory_panels} shows representative test trajectories from
each system alongside the corresponding observer score sequences. The panels
display how scores rise in the pre-bifurcation region for signal trajectories
and stay flat for null trajectories. A method with high detection time should
show early score rises in these panels; a method with low false-positive rate
should produce near-zero scores on null sequences.

\subsection{Training Behavior}
\label{sec:design:training}

Figure~\ref{fig:training_curves} plots training and validation loss curves for
both variants at the 500-patient setting. These curves confirm that training
converges within the allotted 50~epochs and show whether early stopping
activates at a consistent epoch or varies across seeds.

\subsection{Supplementary and Exploratory Comparisons}
\label{sec:design:supplementary}

Several comparisons are excluded from the main tables. Classical early-warning
indicators (lag-1 autocorrelation, rolling variance) were evaluated during
development but use a different scoring pipeline and are not included in the
finalized benchmark. Ablation experiments---EWS-augmented input features,
auxiliary spectral heads, null-training configurations---were run with fewer
seeds or on a subset of systems. These results are retained in the experiment
directory for reference but are not mixed into the main claims, which report
only the two observer variants across the complete seed and system grid.

\section{Results}
\label{sec:results}

\subsection{No Single Observer Dominates Across Systems and Metrics}
\label{sec:results:overview}

Table~\ref{tab:benchmark_500} summarizes the 500-patient benchmark.

On the fold system, Kalman-LSTM-Spec achieves higher early-warning AUC
(\(0.903 \pm 0.059\) vs.\ \(0.827 \pm 0.235\)) with lower cross-seed variance.
Kalman-BCE produces higher mean detection time (\(109.5 \pm 28.6\) vs.\
\(72.7 \pm 11.0\) steps), indicating earlier threshold crossings, and a
slightly lower false-positive rate (\(0.330\) vs.\ \(0.379\)).

On the Hopf system, both methods maintain low false-positive rates
(\(\approx 0.05\)). LSTM-Spec again has higher EW-AUC (\(0.823 \pm 0.243\)
vs.\ \(0.617 \pm 0.177\)). BCE has higher mean detection time (\(60.3\) vs.\
\(49.5\) steps) but with larger variance (\(\pm 26.5\) vs.\ \(\pm 3.8\)).

On the logistic system, both methods reach perfect EW-AUC (\(1.000\)) and
similar detection times (\(\approx 62\) steps). The only separation is
false-positive rate: BCE is lower (\(0.201 \pm 0.143\) vs.\
\(0.268 \pm 0.098\)).

No method dominates all three metrics on any system. LSTM-Spec tends to
discriminate better (higher EW-AUC on fold and Hopf), while BCE is preferable
when false-positive control is the priority or when earlier raw detection time
matters.

\subsection{Patient Count Changes the Method Ranking}
\label{sec:results:sweep}

Table~\ref{tab:benchmark_sweep} and Figure~\ref{fig:patient_sweep} report
performance across patient counts from 100 to 500.

On the fold system, LSTM-Spec detection time decreases from \(108\) to \(73\)
steps as patient count grows, while BCE detection time stays near
\(110\)--\(121\). The EW-AUC gap widens in the same direction: both methods
start near \(0.87\)--\(0.89\) at 100~patients, but by 500~patients LSTM-Spec
reaches \(0.903\) while BCE is at \(0.827\). LSTM-Spec benefits more from
additional training data on this system.

On the Hopf system, LSTM-Spec EW-AUC increases from \(0.652\) to \(0.823\)
across the sweep, compared with a flat range of \(0.60\)--\(0.64\) for BCE.
However, LSTM-Spec detection time is undefined at 200~patients and is reported
from a single seed at 300~patients. The threshold (\(0.500\)) exceeds the
model's pre-bifurcation scores at these patient counts, even though the scores
still discriminate in the EW-AUC sense. By 400--500~patients, threshold
crossings become consistent.

On the logistic system, both methods reach EW-AUC~\(= 1.0\) by 200~patients
and maintain stable detection times near \(62\)--\(66\) steps. False-positive
rate increases with patient count for both methods, with LSTM-Spec consistently
higher.

\subsection{Score Trajectories Reveal Regime-Specific Behavior}
\label{sec:results:trajectories}

Figure~\ref{fig:trajectory_panels} shows representative score trajectories at
500~patients.

On the fold system, the BCE score rises gradually and crosses threshold early in
the trajectory. The LSTM-Spec score holds low longer and then rises sharply near
the bifurcation. This matches the aggregate pattern: BCE has higher mean
detection time (earlier crossings) but wider variance, while LSTM-Spec has a
tighter, later crossing with better signal--null separation.

On the Hopf system, the BCE score is nearly flat. The LSTM-Spec score shows more
structure near the bifurcation but, in the displayed seed, does not cross
threshold before the transition.

On the logistic system, the BCE score remains below threshold for the displayed
seed. The LSTM-Spec score crosses threshold early and drops after the
bifurcation point.

\subsection{Training Curves and Optimization Behavior}
\label{sec:results:training}

Figure~\ref{fig:training_curves} plots training and validation loss for both
variants at 500~patients. Both methods converge smoothly with minimal
train--validation gap. BCE plateaus around epoch~20--30, while LSTM-Spec
continues to decrease through the full 50~epochs. LSTM-Spec reaches a lower
final loss on all three systems.

The lower training loss for LSTM-Spec is consistent with its higher EW-AUC on
fold and Hopf: the model fits the Gaussian ramp target more tightly. On the
logistic system, both methods reach perfect EW-AUC despite a large loss
difference, indicating that the logistic transition is detectable even with a
coarser fit to the training target.

\subsection{Summary of Main Findings}
\label{sec:results:summary}

\begin{itemize}[leftmargin=*]
  \item LSTM-Spec achieves higher EW-AUC than BCE on fold and Hopf at
        500~patients. Both reach perfect EW-AUC on logistic.
  \item BCE produces earlier detections (higher DT) on fold and comparable DT
        elsewhere, but with larger cross-seed variance.
  \item BCE has lower false-positive rate on fold and logistic. Both methods
        have low FPR on Hopf at high patient counts.
  \item LSTM-Spec benefits more from increasing patient count on fold and Hopf.
        Logistic performance saturates by 200~patients for both methods.
  \item Score trajectories differ in shape: BCE rises gradually, LSTM-Spec
        produces sharper transitions closer to the bifurcation.
\end{itemize}

\section{Discussion}
\label{sec:discussion}

% Interpret the results and explain why they matter.

\subsection{Regime Dependence Is the Main Finding}
\label{sec:discussion:regime}

% Different bifurcation mechanisms create different detection requirements.

\subsection{Why Spectral or Recurrent Structure Helps Only in Some Settings}
\label{sec:discussion:spectral}

% Inductive bias discussion:
% \begin{itemize}
%   \item Useful when the dynamics match it.
%   \item Not expected to dominate every transition type.
%   \item Not a failure if it loses where the bias is mismatched.
% \end{itemize}

\subsection{Why Simple Baselines Remain Important}
\label{sec:discussion:baselines}

% Simple indicators can be strong because they encode direct dynamical signatures.
% This strengthens the paper because it shows the comparison is not weak.

\subsection{Practical Implications}
\label{sec:discussion:implications}

% Observer choice should depend on:
% \begin{itemize}
%   \item Dynamical regime
%   \item Tolerance for false positives
%   \item Desired warning lead time
%   \item Available data depth
% \end{itemize}

\section{Limitations}
\label{sec:limitations}

% Show reviewers that the study is honest and bounded.
%
% Points to cover:
% \begin{itemize}
%   \item Synthetic systems only.
%   \item Limited set of observer variants in the finalized main benchmark.
%   \item Threshold calibration may affect detection-time conclusions.
%   \item Patient-count sweeps do not cover all possible data regimes.
%   \item Results should not be generalized to all domains without additional
%         validation.
% \end{itemize}
%
% Do not hide limitations. Use them to make the claims precise.

\section{Conclusion}
\label{sec:conclusion}

% Close with the central scientific message.
%
% Learned Kalman observers show useful but conditional early-warning behavior.
% Across fold, Hopf, and logistic systems, the preferred observer depends on
% the dynamical regime, the metric, and the available data. The study therefore
% supports regime-aware observer design rather than a universal-detector
% interpretation.

% --- Back Matter ---

\section*{Data Availability}
% Statement about where the generated data is available.

\section*{Code Availability}
% Statement about where the code is available.

\section*{Acknowledgements}
% Funding sources, colleagues, etc.

\section*{Author Contributions}
% Who did what.

\section*{Conflicts of Interest}
% Declaration of competing interests.


\appendix
\section{Supplementary Material}
\label{sec:appendix}

\subsection{System Equations and Parameters}
\label{sec:appendix:systems}

Table~\ref{tab:system_params} lists the parameter defaults for each system.
Bifurcation time is computed from the control-parameter sweep so that the
bifurcation occurs at a known time step within each trajectory.

\begin{table}[htbp]
  \centering
  \caption{Data-generation parameters for each dynamical system. All
  trajectories have length $T = 200$. The bifurcation time $\tau$ is derived
  from the control-parameter endpoints and the bifurcation value.}
  \label{tab:system_params}
  \small
  \begin{tabular}{llll}
    \toprule
    Parameter & Fold & Hopf & Logistic \\
    \midrule
    Update rule
      & $x \leftarrow x + (r - x^2) + \xi$
      & $\dot{r} = \mu r - r^3 + \xi$
      & $x \leftarrow \mu x(1 - x) + \xi$ \\[2pt]
    Control parameter
      & $r \in [2.0, -1.0]$
      & $\mu \in [-0.5, 0.5]$
      & $\mu \in [2.5, 4.0]$ \\[2pt]
    Bifurcation value
      & $r^{*} = 0$
      & $\mu^{*} = 0$
      & $\mu^{*} = 3.0$ \\[2pt]
    $\tau$ (steps)
      & $133$
      & $100$
      & $67$ \\[2pt]
    Process noise $\sigma_{\xi}$
      & 0.30
      & 0.05
      & 0.02 \\[2pt]
    Observation noise $\sigma_{\mathrm{obs}}$
      & 0.10
      & 0.15
      & 0.05 \\[2pt]
    Observation dim $D$
      & 1
      & 2
      & 1 \\[2pt]
    Null behavior
      & $r$ fixed at $r_{\mathrm{start}}$
      & $\mu$ fixed at $\mu_{\mathrm{start}}$
      & $\mu$ fixed at 2.8 \\
    \bottomrule
  \end{tabular}
\end{table}

\textbf{Fold.} The state $x$ follows $x_{t+1} = x_t + (r_t - x_t^2) +
\sigma_\xi\,\varepsilon_t$ with $\varepsilon_t \sim \mathcal{N}(0,1)$. The
control parameter $r$ decreases linearly from 2.0 to $-1.0$ over 200~steps.
At $r = 0$ the stable equilibrium is lost. Null trajectories hold $r$ at
$r_{\mathrm{start}} = 2.0$ for all steps. Initial state:
$x_0 \sim \mathcal{N}(0, 0.5^2)$. State clipped to $[-5, 5]$.

\textbf{Hopf.} The system is written in polar coordinates $(r, \theta)$:
$r_{t+1} = r_t + (\mu_t\,r_t - r_t^3) + \sigma_\xi\,\varepsilon_t$ and
$\theta_{t+1} = \theta_t + \omega + \sigma_\xi\,\varepsilon_t'$, with
$\omega = 0.1$. The observation is $(r\cos\theta, r\sin\theta)$ plus
observation noise. The control parameter $\mu$ increases from $-0.5$ to $0.5$;
the bifurcation occurs at $\mu = 0$. Initial conditions:
$r_0 \sim U(0.5, 1.5)$, $\theta_0 \sim U(0, 2\pi)$. Null trajectories hold
$\mu$ at $\mu_{\mathrm{start}} = -0.5$.

\textbf{Logistic.} The map $x_{t+1} = \mu_t\,x_t(1 - x_t) +
\sigma_\xi\,\varepsilon_t$ with $\mu$ swept from 2.5 to 4.0. The first
period-doubling bifurcation occurs at $\mu = 3.0$. Initial state:
$x_0 \sim U(0.1, 0.9)$. Null trajectories hold $\mu$ at 2.8.

\subsection{Training and Model Hyperparameters}
\label{sec:appendix:hyperparams}

Table~\ref{tab:hyperparams} lists all training and model hyperparameters used
in the finalized benchmark.

\begin{table}[htbp]
  \centering
  \caption{Training and model hyperparameters. Values are fixed across all
  systems and seeds.}
  \label{tab:hyperparams}
  \small
  \begin{tabular}{ll}
    \toprule
    Hyperparameter & Value \\
    \midrule
    \multicolumn{2}{l}{\emph{Model}} \\[2pt]
    Latent dimension $d$ & 4 \\
    LSTM hidden dimension (LSTM-Spec only) & 8 \\
    Transition init $A$ & $0.95\,I_d$ \\
    $C$, $K$ init & $\mathcal{N}(0, 0.1^2)$ \\[4pt]
    \multicolumn{2}{l}{\emph{Training}} \\[2pt]
    Optimizer & AdamW \\
    Learning rate & $10^{-3}$ \\
    Weight decay & $10^{-5}$ \\
    Scheduler & Cosine annealing \\
    Scheduler $\eta_{\min}$ & $10^{-6}$ \\
    Epochs & 50 \\
    Batch size & 64 \\
    Gradient clip norm & 1.0 \\
    Early-stopping patience & 5 \\[4pt]
    \multicolumn{2}{l}{\emph{Loss}} \\[2pt]
    Primary loss & Binary cross-entropy \\
    Gaussian ramp $\sigma$ & 60 \\
    Spectral penalty weight $\lambda$ (LSTM-Spec) & 0.01 \\
    Spectral threshold & 0.95 \\
    Power iterations for $\rho(M)$ & 10 \\[4pt]
    \multicolumn{2}{l}{\emph{Data split}} \\[2pt]
    Train / Validation / Test & 60\% / 20\% / 20\% \\
    \bottomrule
  \end{tabular}
\end{table}

\subsection{Batch Inventory}
\label{sec:appendix:batches}

Table~\ref{tab:benchmark_inventory} lists all experiment batches in the
results directory. For each patient count, the batch-selection rule
(Section~\ref{sec:batch}) selects the batch with the most complete coverage.
Partial batches from development runs are retained but excluded from the
reported figures and tables.

\begin{table}[htbp]
  \centering
  \caption{Inventory of benchmark batches. The selected batch for each patient
  count is the most complete batch available.}
  \label{tab:benchmark_inventory}
  \begin{tabular}{llrrrl}
    \toprule
    Patients & Batch & Records & Unique triples & Systems & Selected \\
    \midrule
    100 & 2026-07-16\_08-11-32 & 60 & 60 & 3 & yes \\
    200 & 2026-07-16\_08-11-59 & 60 & 60 & 3 & yes \\
    300 & 2026-07-16\_08-08-14 & 6 & 6 & 1 & no \\
    300 & 2026-07-16\_08-12-57 & 60 & 60 & 3 & yes \\
    400 & 2026-07-16\_08-08-39 & 3 & 3 & 1 & no \\
    400 & 2026-07-16\_08-13-29 & 60 & 60 & 3 & yes \\
    500 & 2026-07-16\_08-13-49 & 60 & 60 & 3 & yes \\
    \bottomrule
  \end{tabular}
\end{table}


\clearpage
\bibliographystyle{plainnat}
\bibliography{bib/references}

\end{document}
